% Part: set-theory
% Chapter: z
% Section: infinity-again

\documentclass[../../../include/open-logic-section]{subfiles}

\begin{document}

\olfileid{sth}{z}{infinity-again}
\olsection{অসীমতা}

এ পর্যন্ত পাওয়া স্বতঃসিদ্ধগুলি থেকে নিশ্চিত করা যায় যে, কোনো সেট
থাকলে অসীমসংখ্যক সেট আছে। কারণ, ধরি কোনো সেট আছে। তাহলে
\olref[sth][z][sep]{prop:emptyexists} অনুযায়ী $\emptyset$-ও আছে।
এখন যেকোনো সেট~$x$-এর জন্য
\olref[sth][z][pairs]{prop:pairsconsequences} অনুযায়ী
$x \cup \{x\}$ সেটটিও আছে। এই ধাপ কয়েকবার প্রয়োগ করলে
নিচের সেটগুলি পাই:
\begin{enumerate}
	\item[0.] $\emptyset$ %& $0$ & stage $0$\\
	\item[1.] $ \{\emptyset\}$ %& $1$ & stage $1$\\
	\item[2.] $\{\emptyset, \{\emptyset\}\}$ %& $2$ & stage $2$\\
	\item[3.] $\{\emptyset, \{\emptyset\}, \{\emptyset, \{\emptyset\}\}\}$ %&$3$ & stage $3$\\
	\item[4.] $\{\emptyset, \{\emptyset\}, \{\emptyset, \{\emptyset\}\}, \{\emptyset, \{\emptyset\}, \{\emptyset, \{\emptyset\}\}\}\}$% & $4$ & stage $4$
\end{enumerate}
যাচাই করা যায়, সেটগুলির প্রত্যেকটি অন্যগুলি থেকে আলাদা।

আমরা গণনা $0$ থেকে শুরু করেছি, তার কয়েকটি কারণ আছে। একটি
কারণ এই: যে সেটটির পাশে ``$n$'' লিখেছি, তাতে ঠিক $n$টি সদস্য
আছে এবং (স্বজ্ঞাগতভাবে) সেটটি $n$-তম পর্যায়ে গঠিত হয়—এটি
যাচাই করা কঠিন নয়।

কিন্তু এখান থেকে
\emph{অসীমসংখ্যক} সেট পাওয়া গেলেও কোনো
\emph{অসীম সেট}, অর্থাৎ অসীমসংখ্যক সদস্যবিশিষ্ট একটি সেট,
থাকার নিশ্চয়তা মেলে না। এই পার্থক্যটি জরুরি: কোনো
(ডেডেকিন্ড-)অসীম সেট না পেলে আমরা ডেডেকিন্ড বীজগঠন
নির্মাণ করতে পারব না। অথচ স্বাভাবিক সংখ্যার সেট হিসেবে
ব্যবহার করার জন্য আমরা এমন একটি বীজগঠন চাই।
(তুলনা করো \olref[sfr][infinite][dedekindsproof]{sec}।)

এ পর্যন্ত স্থির করা স্বতঃসিদ্ধগুলি যে কোনো অসীম সেটের
অস্তিত্ব
\emph{নিশ্চিত করে না}, সেটি মনে রাখা দরকার। তাই
একটি নতুন স্বতঃসিদ্ধ নিতে হবে:

\begin{axiom}[অসীমতা]
এমন একটি সেট $I$ আছে, যাতে $\emptyset \in I$ এবং
যখনই $x \in I$, তখনই $x \cup \{x\} \in I$।
\begin{align*}
	\exists I( & (\exists o \in I)\forall x\ x \notin o \land {}\\
	& (\forall x \in I)(\exists s \in I)\forall z(z \in s \liff (z \in x \lor z = x)))
\end{align*}
\end{axiom}

অসীমতার স্বতঃসিদ্ধে পাওয়া সেট $I$ যে ডেডেকিন্ড-অসীম,
তা সহজেই দেখা যায়। এর বিশিষ্ট সদস্য $\emptyset$; আর $I$-এর
উপর এক-এক চিত্রণটি $s(x) = x\cup \{x\}$ দিয়ে নির্ধারিত।
এখন \olref[sfr][infinite][dedekind]{thm:DedekindInfiniteAlgebra}-এ
দেখানো হয়েছে কীভাবে একটি ডেডেকিন্ড-অসীম সেট থেকে ডেডেকিন্ড
বীজগঠন পাওয়া যায়। আমরা সেটিকেই স্বাভাবিক সংখ্যার সেট
হিসেবে ধরব। আরও নির্দিষ্টভাবে:
\begin{defn}\ollabel{defnomega}
ধরি, $I$ অসীমতার স্বতঃসিদ্ধে পাওয়া যেকোনো সেট। ধরি,
$s$ হলো $s(x) = x \cup \{x\}$ অপেক্ষকটি। ধরি,
$\omega = \closureofunder{s}{\emptyset}$। $\omega$-এর
সদস্যদের আমরা \emph{স্বাভাবিক সংখ্যা} বলব, এবং বলব যে,
$\emptyset$-এর উপর $s$-কে $n$-বার প্রয়োগের ফল হলো $n$।
\end{defn}

এখন কয়েক অনুচ্ছেদ আগের ``$n$'' চিহ্নিত সেটটি আবার দেখলে
যাচাই করতে পারবে যে, সেটিকেই সংখ্যা $n$
\emph{হিসেবে} ধরা হবে।

এই নির্ধারণের তাৎপর্য নিয়ে
\olref[sth][z][nat]{sec}-এ আলোচনা করব। আপাতত এর সাহায্যে
একটি স্বজ্ঞাসম্মত ফল প্রমাণ করা যায়:

\begin{prop}\ollabel{naturalnumbersarentinfinite}
কোনো স্বাভাবিক সংখ্যাই ডেডেকিন্ড-অসীম নয়।
\end{prop}

\begin{proof}
প্রমাণটি আবেশের সাহায্যে, অর্থাৎ
\olref[sfr][infinite][induction]{thm:dedinfiniteinduction} অনুযায়ী।
স্পষ্টতই $0 = \emptyset$ ডেডেকিন্ড-অসীম নয়। আবেশের ধাপে
আমরা বিপরীতার্থক দাবিটি প্রমাণ করব: যদি (অসঙ্গতভাবে)
$s(n)$ ডেডেকিন্ড-অসীম হয়, তবে $n$-ও ডেডেকিন্ড-অসীম।

তাই ধরি, $s(n)$ ডেডেকিন্ড-অসীম। অর্থাৎ এমন একটি
!!{injection} $f$ আছে, যার জন্য
$\ran{f}\subsetneq \dom{f} = s(n) = n \cup \{n\}$।
এখন দুটি ক্ষেত্র আছে।

\emph{প্রথম ক্ষেত্র: $n \notin \ran{f}$।} তাহলে
$\ran{f} \subseteq n$ এবং $f(n) \in n$। ধরি,
$g = \funrestrictionto{f}{n}$; তাহলে
$\ran{g} = \ran{f} \setminus \{f(n)\} \subsetneq n = \dom{g}$।
সুতরাং $n$ ডেডেকিন্ড-অসীম।

\emph{দ্বিতীয় ক্ষেত্র: $n \in \ran{f}$।} একটি
$m \in \dom{f} \setminus \ran{f}$ নিই এবং সংজ্ঞাক্ষেত্র
$s(n) = n \cup \{n\}$-বিশিষ্ট একটি অপেক্ষক $h$ এভাবে
সংজ্ঞায়িত করি:
\[
h(x) =
	\begin{cases}
		f(x) & \text{যদি }f(x) \neq n\\
		m & \text{যদি }f(x)=n
	\end{cases}
\]
অতএব $h$ ও $f$ সর্বত্র এক, কেবল
$h(f^{-1}(n)) = m \neq n = f(f^{-1}(n))$। যেহেতু
$f$ একটি !!a{injection}, তাই $n \notin \ran{h}$;
আর $\ran{h} \subsetneq \dom{h} = s(n)$। এবার প্রথম
ক্ষেত্রের যুক্তি প্রয়োগ করে দেখা যায়, $n$ ডেডেকিন্ড-অসীম।
\end{proof}

তবে অসীমতার স্বতঃসিদ্ধকে কীভাবে
\emph{ন্যায্যতা} দেব, সেই প্রশ্ন রয়ে যায়। সংক্ষিপ্ত উত্তর
হলো, এতক্ষণ পর্যায়গুলি নিয়ে যে বর্ণনা দিয়েছি তাতে আরেকটি
নীতি যোগ করতে হবে। নীতিটি এই:
\begin{enumerate}
	\item[] \stagesinf. একটি অসীম পর্যায় আছে। অর্থাৎ এমন একটি পর্যায় আছে যা (ক) প্রথম পর্যায় নয়, (খ) যার আগে কিছু পর্যায় আছে, কিন্তু (গ) যার ঠিক আগের কোনো পর্যায় নেই।
\end{enumerate}
এই নীতি থেকে অসীমতার স্বতঃসিদ্ধ সহজেই পাওয়া যায়। আমরা
জানি, স্বাভাবিক সংখ্যা $n$ গঠিত হয় পর্যায়~$n$-এ। তাই
$\omega$ গঠিত হয় প্রথম অসীম পর্যায়ে। আর $\omega$ নিজেই
অসীমতার স্বতঃসিদ্ধের সাক্ষী।

অবশ্য এতে প্রশ্নটি শুধু পিছিয়ে যায়: \stagesinf{}-কে
ন্যায্যতা দেব কীভাবে? \stagessucc-এর মতো এটিও সেটের
সঞ্চয়ী-পর্যায়ক্রমিক ক্রমধারা নিয়ে আগের বর্ণনার স্পষ্ট অংশ ছিল
না। তার চেয়েও বড় কথা, যে ধারণায় সেট পর্যায়ক্রমে গঠিত
হয়, সেই ধারণা থেকেই গঠনপ্রক্রিয়ায় একটি
\emph{অসীম} পর্যায় থাকতেই হবে বলে সিদ্ধান্ত আসে না।
পর্যায়গুলি স্বাভাবিক সংখ্যার মতো ক্রমানুসারে সাজানো—এমন
ভাবনাও পুরোপুরি সংগত মনে হয়।

কিন্তু এখানেই একটি স্পষ্ট সমস্যা দেখা দেয়।
\olref[sfr][infinite][dedekindsproof]{sec}-এ আমরা
(চিন্তার) একটি ডেডেকিন্ড-অসীম সেট আছে বলে ডেডেকিন্ডের
``প্রমাণ'' দেখেছিলাম। সেটি খুব সন্তোষজনক না-ও মনে হতে
পারে। কিন্তু সেটের ধাপে-ধাপে গঠনের ধারণা (বা ``চিন্তার
নিয়ম'') যদি \stagesinf{}-কে আমাদের উপর ``চাপিয়ে না দেয়'',
তবে ডেডেকিন্ড-অসীম সেট আছে—এই দাবির কোনো অন্তর্নিহিত
ন্যায্যতা এখনও আমাদের হাতে নেই।

এ বিষয়ে আরও অনেক কিছু বলার আছে। তবে আশা করি, এখন
একটি স্বতঃসিদ্ধ (বা নীতি) ন্যায্যতা পেতে কী
\emph{প্রয়োজন} হতে পারে, তা নিয়ে ভাবতে পারবে। এর পরের
আলোচনায় আমরা \stagesinf{}-কে স্বীকার করে নেব।

\end{document}
